\begin{abstract}
\label{sec:abstract}

Long-horizon scientific agents produce hundreds of tool invocations, state
transitions, and intermediate claims before arriving at a terminal result.
Final-task accuracy obscures when, why, and whether an error became
unrecoverable. This paper introduces \systemname, a governed benchmark
and measurement framework that makes failure dynamics—onset, propagation,
and closing recovery windows—empirically observable across scientific-agent
trajectories. We specify a causal fault model with seven mechanism families,
four workload levels spanning repository repair to nonstationary-noise
experiments, held-out task, fault, and composition splits, deterministic
fault injectors with clean counterfactuals, and an independent adjudication
protocol that separates execution correctness from scientific validity. A
calibrated early-warning and attribution baseline, built from
execution-grounded state-difference histories and evidence provenance,
is evaluated on the frozen benchmark to establish whether pre-terminal
failure prediction transfers to unseen task and fault families. The full
artifact—trace schema, manifests, injectors, adjudicator protocol, baseline
code, and evaluation harness—is released under a documented access,
licensing, and maintenance plan.

\end{abstract}
